% Wersja źródłowa: audyt II, jednostka 20
\setcounter{section}{19}
\section{Pochodne dystrybucyjne i transpozycja operatorów}
\label{sec:pochodne20}

\subsection{Definicja pochodnej i ciągłość w topologiach dualnych}

W tym rozdziale $d\ge1$ jest liczbą całkowitą,
$\Omega\subseteq\R^d$ jest niepustym zbiorem otwartym,
a $\K\in\{\R,\C\}$. Dystrybucje są
ciągłymi $\K$-liniowymi funkcjonałami na $\mathcal D(\Omega;\K)$, zgodnie
z konwencją przyjętą w \hyperref[sec:dystrybucje19]{rozdziale~19}.
Dla wielowskaźnika $\alpha=(\alpha_1,\ldots,\alpha_d)\in\N_0^d$
piszemy $|\alpha|=\sum_j\alpha_j$,
$\partial^\alpha=\partial_1^{\alpha_1}\cdots\partial_d^{\alpha_d}$
oraz $\binom\alpha\beta=\prod_j\binom{\alpha_j}{\beta_j}$
dla $\beta\le\alpha$ współrzędnie. Przez $e_j$ oznaczamy
$j$-ty wektor bazy standardowej $\R^d$, a przez $\mathcal D_K$
i $p_m$ przestrzenie i seminormy z rozdziału~18.
Pominięcie $\K$ w oznaczeniu przestrzeni funkcji lub dystrybucji
oznacza, że nadal pracujemy nad ustalonym ciałem $\K$.

\begin{lemat}[Różniczkowanie funkcji testowych]
\label{lem:test-derivative20}
Dla każdego $\alpha\in\N_0^d$ odwzorowanie
$D_\alpha:\mathcal D(\Omega;\K)\to\mathcal D(\Omega;\K)$,
$D_\alpha\varphi=\partial^\alpha\varphi$, jest ciągłe i liniowe.
Ponadto dla każdego zbioru ograniczonego $B\subseteq\mathcal D(\Omega)$
zbiór $D_\alpha(B)$ jest ograniczony.
\end{lemat}
\begin{proof}
Dla każdego $K\Subset\Omega$ odwzorowanie $D_\alpha$ zachowuje
$\mathcal D_K(\Omega)$, ponieważ w otoczeniu każdego punktu
spoza $\supp\varphi$ funkcja $\varphi$ jest identycznie zerowa,
a więc wszystkie jej pochodne również są tam zerowe.
Przy seminormach z \hyperref[sec:funkcje-testowe18]{rozdziału~18} mamy
$p_m(D_\alpha\varphi)\le p_{m+|\alpha|}(\varphi)$.
Ograniczenie $D_\alpha|_{\mathcal D_K}$ jest więc ciągłe.
Po złożeniu z ciągłym włożeniem $\mathcal D_K\to\mathcal D(\Omega)$
otrzymujemy ciągłość na każdym etapie, a z
Twierdzenia~\ref{thm:LF-universal18} ciągłość na całej $\mathcal D(\Omega)$.
Liniowość wynika z liniowości różniczkowania.
Jeśli $B$ jest ograniczony i $U$ jest otoczeniem zera
w $\mathcal D(\Omega)$, to $D_\alpha^{-1}(U)$ jest
otoczeniem zera. Istnieje $t>0$ takie, że
$B\subseteq tD_\alpha^{-1}(U)$, a z liniowości
$D_\alpha(B)\subseteq tU$. Zatem $D_\alpha(B)$ jest ograniczony.
\end{proof}

\begin{definicja}[Pochodna dystrybucyjna]
\label{def:dist-deriv20}
Niech $T\in\mathcal D'(\Omega;\K)$ i $\alpha\in\N_0^d$.
\emph{Pochodną dystrybucyjną} $\partial^\alpha T$ określamy wzorem
\[
 \langle\partial^\alpha T,\varphi\rangle
 =(-1)^{|\alpha|}\langle T,\partial^\alpha\varphi\rangle,
 \qquad\varphi\in\mathcal D(\Omega;\K).
\]
Prawa strona jest ciągłym funkcjonałem liniowym względem $\varphi$,
na mocy Lematu~\ref{lem:test-derivative20}.
Dla $j\in\{1,\ldots,d\}$ i $\alpha=e_j$ piszemy również $\partial_jT$.
\end{definicja}

\begin{twierdzenie}[Ciągłość i algebra pochodnych]
\label{thm:dist-deriv-cont20}
Niech $T\in\mathcal D'(\Omega;\K)$ oraz
$\alpha,\beta\in\N_0^d$. Wówczas:
\begin{enumerate}
\item $\partial^\alpha T$ jest dystrybucją dla każdego
$T\in\mathcal D'(\Omega)$;
\item $\partial^\alpha\partial^\beta T
=\partial^{\alpha+\beta}T$, w szczególności
$\partial_j\partial_kT=\partial_k\partial_jT$
dla $j,k\in\{1,\ldots,d\}$;
\item operator $\partial^\alpha:\mathcal D'(\Omega)\to
\mathcal D'(\Omega)$ jest ciągły zarówno w topologii słabej-*
$\sigma(\mathcal D',\mathcal D)$, jak i silnej
$\beta(\mathcal D',\mathcal D)$;
\item $\supp(\partial^\alpha T)\subseteq\supp T$.
Jeśli dla zwartego $K\Subset\Omega$ i $m\in\N_0$
dystrybucja $T$ ma na $K$ rząd nie większy niż $m$,
to $\partial^\alpha T$ ma na $K$ rząd nie większy niż
$m+|\alpha|$.
\end{enumerate}
\end{twierdzenie}
\begin{proof}
Punkt pierwszy wynika z ciągłości $D_\alpha$ na
$\mathcal D(\Omega)$ i złożenia funkcjonałów.
Dla każdej $\varphi\in\mathcal D(\Omega)$ dwukrotne
zastosowanie definicji daje
\[
\begin{aligned}
\langle\partial^\alpha\partial^\beta T,\varphi\rangle
&=(-1)^{|\alpha|+|\beta|}
  \langle T,\partial^\beta\partial^\alpha\varphi\rangle\\
&=(-1)^{|\alpha+\beta|}
  \langle T,\partial^{\alpha+\beta}\varphi\rangle
 =\langle\partial^{\alpha+\beta}T,\varphi\rangle.
\end{aligned}
\]
W drugiej równości korzystamy z przemienności pochodnych
funkcji gładkich i z $|\alpha+\beta|=|\alpha|+|\beta|$.
W szczególności można zamienić kolejność dwóch pochodnych
pierwszego rzędu. Operator $\partial^\alpha$ jest liniowy
z definicji.
Dla ustalonej $\varphi$ zachodzi
$q_\varphi(\partial^\alpha T)=q_{\partial^\alpha\varphi}(T)$,
więc otrzymujemy ciągłość słabą-*.
Dla zbioru ograniczonego $B\subseteq\mathcal D(\Omega)$
mamy natomiast
\[
 q_B(\partial^\alpha T)=q_{D_\alpha(B)}(T).
\]
Z lematu~\ref{lem:test-derivative20} zbiór $D_\alpha(B)$ jest
ograniczony, co daje ciągłość silną. Gdy $T$ zeruje się na
otwartym $U\subseteq\Omega$, zeruje się tam także $\partial^\alpha T$.
Istotnie, jeśli $e_U$ oznacza przedłużenie zerem z definicji
\ref{def:distribution-zero19}, to dla $\varphi\in\mathcal D(U)$
mamy $\partial^\alpha(e_U\varphi)=e_U(\partial^\alpha\varphi)$.
Równość zachodzi na $U$ i w otoczeniu każdego punktu
spoza $U$, gdzie obie strony znikają dzięki zwartemu nośnikowi
$\varphi$ w $U$. Stąd
$\langle(\partial^\alpha T)|_U,\varphi\rangle
=(-1)^{|\alpha|}\langle T|_U,\partial^\alpha\varphi\rangle=0$.
Stosując to do $U=\Omega\setminus\supp T$, otrzymujemy
inkluzję nośników.
Wreszcie dla $\varphi\in\mathcal D_K$ mamy
$|\langle\partial^\alpha T,\varphi\rangle|
\le C_Kp_m(\partial^\alpha\varphi)
\le C_Kp_{m+|\alpha|}(\varphi)$.
\end{proof}

\begin{stwierdzenie}[Zgodność z pochodną klasyczną]
\label{prop:classical-weak20}
Niech $m\in\N$, $f\in C^m(\Omega;\K)$, a $\alpha\in\N_0^d$
niech spełnia $|\alpha|\le m$. Wtedy
$\partial^\alpha T_f=T_{\partial^\alpha f}$.
Dla $f,g\in L^1_{\mathrm{loc}}(\Omega;\K)$
oraz $j\in\{1,\ldots,d\}$ tożsamość
$\partial_jT_f=T_g$ jest równoważna temu, że $g$ jest
\emph{słabą pochodną} $f$ w kierunku $x_j$, tj.
\[
 \int_\Omega f(x)\,\partial_j\varphi(x)\,\mathrm d\lambda_d(x)
 =-\int_\Omega g(x)\varphi(x)\,\mathrm d\lambda_d(x)
 \quad(\varphi\in\mathcal D(\Omega)).
\]
\end{stwierdzenie}
\begin{proof}
Ustalmy $\varphi\in\mathcal D(\Omega)$ i wybierzmy
$\chi\in C_c^\infty(\Omega;\R)$ równą $1$ w otoczeniu
$\supp\varphi$ (Lemat~\ref{lem:cutoffs17}).
Przedłużenie $F$ funkcji $\chi f$ zerem na $\R^d$
jest klasy $C^m$ i ma zwarty nośnik. Niech $\Phi$ będzie
przedłużeniem $\varphi$ zerem na $\R^d$.
Na otoczeniu $\supp\varphi$ funkcje $F$ i $f$ oraz ich pochodne
do rzędu $m$ są zgodne. Wszystkie pochodne $\Phi$
znikają poza $\supp\varphi$.
Dla $j\in\{1,\ldots,d\}$ i funkcji $F$ klasy co najmniej $C^1$
twierdzenie Fubiniego i całkowanie przez części względem $x_j$
dają
\[
\int_{\R^d}F\,\partial_j\Phi\,\mathrm d\lambda_d
=-\int_{\R^d}(\partial_jF)\Phi\,\mathrm d\lambda_d.
\]
Możemy całkować po dostatecznie dużym prostokącie zawierającym
oba nośniki we wnętrzu, więc wyrazy brzegowe znikają.
Całki są bezwzględnie zbieżne dzięki ciągłości i zwartym nośnikom.
Powtarzając argument dla kolejnych pochodnych zgodnie z
wielowskaźnikiem $\alpha$, dostajemy
\[
\begin{aligned}
\langle\partial^\alpha T_f,\varphi\rangle
&=(-1)^{|\alpha|}\int_{\R^d}F\partial^\alpha\Phi\,\mathrm d\lambda_d\\
&=\int_{\R^d}(\partial^\alpha F)\Phi\,\mathrm d\lambda_d
=\int_\Omega(\partial^\alpha f)\varphi\,\mathrm d\lambda_d
=\langle T_{\partial^\alpha f},\varphi\rangle.
\end{aligned}
\]
Dla $\alpha=0$ jest to tożsamość bez całkowania przez części.
Nie potrzeba żadnych warunków brzegowych na $f$.

W drugiej części z definicji
$\langle\partial_jT_f,\varphi\rangle
=-\int_\Omega f\partial_j\varphi\,\mathrm d\lambda_d$,
natomiast $\langle T_g,\varphi\rangle
=\int_\Omega g\varphi\,\mathrm d\lambda_d$.
Równość dystrybucji zachodzi dokładnie wtedy, gdy te wartości
są równe dla wszystkich funkcji testowych.
Jeżeli słaba pochodna istnieje, jest jednoznaczna jako element
$L^1_{\mathrm{loc}}(\Omega;\K)$ na mocy
Stwierdzenia~\ref{prop:regular-injective19}.
\end{proof}

\subsection{Skoki, funkcja Heaviside'a i osobliwości punktowe}

\begin{przyklad}[Heaviside i delta Diraca]
\label{ex:heaviside20}
Na $\R$ niech $H(x)=0$ dla $x<0$ oraz $H(x)=1$ dla $x>0$;
przyjmijmy $H(0)=0$. Wartość w $0$ nie ma znaczenia
dla wyznaczonej dystrybucji, ponieważ $\lambda_1(\{0\})=0$.
Dla dowolnej $\varphi\in\mathcal D(\R)$ otrzymujemy
\[
 \langle (T_H)',\varphi\rangle
 =-\int_0^\infty\varphi'(x)\,\mathrm d\lambda_1(x)
 =\varphi(0)=\langle\delta_0,\varphi\rangle.
\]
Całka ze znakiem minus jest równa $\varphi(0)$, gdyż $\varphi$ znika
dla dostatecznie dużych $x$.
Zatem $(T_H)'=\delta_0$; w dalszym tekście piszemy krótko
$H'=\delta_0$, utożsamiając funkcje lokalnie całkowalne
z odpowiadającymi im dystrybucjami.
Ponieważ $\operatorname{sgn}x=2H(x)-1$ prawie wszędzie,
a pochodna stałej dystrybucji jest zerowa
(Stwierdzenie~\ref{prop:classical-weak20}), mamy
$(\operatorname{sgn}x)'=2\delta_0$.
Ponadto
\[
\begin{aligned}
\langle(T_{|\cdot|})',\varphi\rangle
&=\int_{-\infty}^0 x\varphi'(x)\,\mathrm d\lambda_1(x)
 -\int_0^\infty x\varphi'(x)\,\mathrm d\lambda_1(x)\\
&=-\int_{-\infty}^0\varphi(x)\,\mathrm d\lambda_1(x)
 +\int_0^\infty\varphi(x)\,\mathrm d\lambda_1(x)
 =\langle T_{\operatorname{sgn}},\varphi\rangle .
\end{aligned}
\]
Wyrazy brzegowe w zerze znikają przez czynnik $x$,
a w nieskończoności przez zwarty nośnik $\varphi$.
Stąd $|x|'=\operatorname{sgn}x$ i $|x|''=2\delta_0$
w sensie dystrybucji.
\end{przyklad}

\begin{center}
\begin{tikzpicture}[scale=0.82,>=Stealth]
\draw[->] (-3,0) -- (3.3,0) node[right] {$x$};
\draw[->] (0,-0.5) -- (0,1.9) node[above] {$H(x)$};
\draw[very thick,blue!60!black] (-2.9,0)--(0,0);
\draw[very thick,blue!60!black] (0,1.3)--(3.0,1.3);
\draw[fill=blue!60!black,draw=blue!60!black,thick] (0,0) circle (2.3pt);
\draw[fill=white,draw=blue!60!black,thick] (0,1.3) circle (2.3pt);
\node[left] at (-0.12,1.3) {$1$};
\node[below left] at (-0.08,-0.05) {$0$};
\draw[gray,dashed] (0,0.12)--(0,1.2);
\end{tikzpicture}

\small\emph{Skok funkcji $H$ z przyjętą wartością $H(0)=0$.
Zmiana wartości w zerze nie wpływa na tożsamość dystrybucyjną
$H'=\delta_0$.}
\end{center}

\begin{twierdzenie}[Pochodna funkcji ze skokiem]
\label{thm:jump20}
Niech $a\in\R$, a $f\in L^1_{\mathrm{loc}}(\R;\K)$ będzie
klasy $C^1$ na $(-\infty,a)$ i $(a,\infty)$. Załóżmy ponadto,
że funkcja ma skończone granice jednostronne
$f(a-),f(a+)$, a klasyczna pochodna określona poza $a$
należy do $L^1_{\mathrm{loc}}(\R;\K)$.
Niech $g(x)=f'(x)$ dla $x\ne a$ (wartość w $a$ jest dowolna).
Wówczas
\[
 (T_f)'=T_g+\bigl(f(a+)-f(a-)\bigr)\delta_a.
\]
\end{twierdzenie}
\begin{proof}
Ustalmy $\varphi\in\mathcal D(\R)$ i $R>0$ takie, że
$\supp\varphi\cup\{a\}\subset(-R,R)$.
Dla $0<\varepsilon<R-|a|$ całkujemy przez części na
$[-R,a-\varepsilon]$ i $[a+\varepsilon,R]$:
\[
\begin{aligned}
-\int_{-R}^{a-\varepsilon}f\varphi'\,\mathrm d\lambda_1
&=-f(a-\varepsilon)\varphi(a-\varepsilon)
 +\int_{-R}^{a-\varepsilon}g\varphi\,\mathrm d\lambda_1,\\
-\int_{a+\varepsilon}^Rf\varphi'\,\mathrm d\lambda_1
&= f(a+\varepsilon)\varphi(a+\varepsilon)
 +\int_{a+\varepsilon}^Rg\varphi\,\mathrm d\lambda_1.
\end{aligned}
\]
Wyrazy brzegowe w $\pm R$ znikają.
Skończone granice jednostronne $f$ zapewniają ograniczoność
$|f|$ przez pewne $M$ w przebitym otoczeniu $a$.
Pominięty wkład do całki z $f\varphi'$ ma moduł
nie większy niż $2\varepsilon M p_1(\varphi)$,
więc dąży do zera.
Ponieważ $g\varphi$ jest całkowalna na $[-R,R]$,
z bezwzględnej ciągłości całki wkład z
$(a-\varepsilon,a+\varepsilon)$ również dąży do zera.
Ciągłość $\varphi$ i granice jednostronne $f$ dają
granicę wyrazów brzegowych
$(f(a+)-f(a-))\varphi(a)$.
Sumując i przechodząc do granicy, otrzymujemy
\[
 -\int_\R f(x)\varphi'(x)\,\mathrm d\lambda_1(x)
 =\int_\R g(x)\varphi(x)\,\mathrm d\lambda_1(x)
 +(f(a+)-f(a-))\varphi(a).
\]
Jest to żądana tożsamość dystrybucyjna.
\end{proof}

\begin{przyklad}[Funkcja częściowo gładka]
\label{ex:piecewise20}
Niech $f(x)=e^x$ dla $x<0$, $f(x)=2e^{-x}$ dla $x>0$
oraz $f(0)=0$. Funkcja jest lokalnie całkowalna,
jest klasy $C^1$ poza zerem, ma granice $f(0-)=1$ i $f(0+)=2$,
a jej pochodna poza zerem jest lokalnie całkowalna.
Z Twierdzenia~\ref{thm:jump20} otrzymujemy, ponieważ
$f(0+)-f(0-)=1$,
\[
 (T_f)'=T_{e^x\chi_{(-\infty,0)}
 -2e^{-x}\chi_{(0,\infty)}}+\delta_0.
\]
Jest to równość dystrybucji na $\R$; jej prawa strona
zawiera składnik osobliwy, którego nie wykrywa klasyczne
różniczkowanie poza punktem skoku.
\end{przyklad}

\begin{przyklad}[Pochodne delty]
\label{ex:delta-derivative20}
Dla $a\in\Omega$ i $\alpha\in\N_0^d$ zachodzi
$\langle\partial^\alpha\delta_a,\varphi\rangle
=(-1)^{|\alpha|}\partial^\alpha\varphi(a)$.
Dla $\Omega=\R$ funkcjonał $\delta_0'$ jest więc różny
od zwykłej miary Diraca: dla funkcji testowej
$\varphi(x)=x\eta(x)$, gdzie $\eta\in C_c^\infty(\R;\R)$
jest równa $1$ w otoczeniu zera, mamy
$\langle\delta_0,\varphi\rangle=0$ oraz
$\langle\delta_0',\varphi\rangle=-1$.
Ze Stwierdzenia~\ref{prop:delta-order19} wynika,
że $\delta_0'$ ma rząd dokładnie $1$ na każdym zwartym
przedziale zawierającym otoczenie zera.
\end{przyklad}

\subsection{Mnożenie przez funkcję gładką i reguła Leibniza}

\begin{definicja}[Mnożenie dystrybucji]
\label{def:multiply20}
Dla $g\in C^\infty(\Omega;\K)$ oraz $T\in\mathcal D'(\Omega;\K)$
definiujemy funkcjonał $gT$ na $\mathcal D(\Omega;\K)$ wzorem
$\langle gT,\varphi\rangle=\langle T,g\varphi\rangle$
dla każdej $\varphi\in\mathcal D(\Omega;\K)$.
Iloczyn $g\varphi$ jest gładki i ma nośnik zawarty w
$\supp\varphi$, więc należy do $\mathcal D(\Omega;\K)$.
Ciągłość funkcjonału $gT$ wykażemy w następnym twierdzeniu.
\end{definicja}

\begin{twierdzenie}[Własności mnożenia]
\label{thm:multiply20}
Niech $g,h\in C^\infty(\Omega;\K)$, $T\in\mathcal D'(\Omega;\K)$
i $\alpha\in\N_0^d$. Wtedy:
\begin{enumerate}
\item $gT$ jest dobrze określoną dystrybucją,
$g(hT)=(gh)T$ oraz $1T=T$;
\item $\supp(gT)\subseteq\supp g\cap\supp T$;
\item operator $M_g:T\mapsto gT$ jest ciągły w topologiach
słabej-* i silnej;
\item zachodzi \emph{wzór Leibniza}
\[
 \partial^\alpha(gT)=
 \sum_{\beta\le\alpha}\binom\alpha\beta
 (\partial^\beta g)\,\partial^{\alpha-\beta}T.
\]
\end{enumerate}
\end{twierdzenie}
\begin{proof}
Oznaczmy przez $N_g$ operator mnożenia przez $g$ na
przestrzeni funkcji testowych, aby odróżnić go od operatora
$M_g$ działającego na dystrybucjach.
Dla każdego niepustego zwartego $K\Subset\Omega$
odwzorowanie $N_g$ przekształca $\mathcal D_K$ w siebie.
Ustalmy $m\in\N_0$. Jeśli $\varphi\in\mathcal D_K$
i $\gamma\in\N_0^d$, $|\gamma|\le m$, to
klasyczna reguła Leibniza daje
$\partial^\gamma(g\varphi)=\sum_{\beta\le\gamma}
\binom\gamma\beta(\partial^\beta g)
(\partial^{\gamma-\beta}\varphi)$.
Wszystkie składniki znikają poza $K$, a na $K$ moduł
$\partial^{\gamma-\beta}\varphi$ nie przekracza $p_m(\varphi)$.
Biorąc supremum i maksimum, otrzymujemy oszacowanie
\[
 p_m(g\varphi)\le C_{K,m,g}\,p_m(\varphi),
 \quad C_{K,m,g}=\max_{|\gamma|\le m}
 \sum_{\beta\le\gamma}\binom\gamma\beta
 \sup_{x\in K}|\partial^\beta g(x)|.
\]
Stała jest skończona ze zwartości $K$ i ciągłości pochodnych $g$.
Dla pustego $K$ etap składa się wyłącznie z funkcji zerowej.
Po złożeniu z włożeniem etapu w $\mathcal D(\Omega)$
otrzymujemy ciągłość $N_g$ na każdym etapie, a z
Twierdzenia~\ref{thm:LF-universal18} ciągłość na $\mathcal D(\Omega)$.
Wobec tego $gT=T\circ N_g$ jest ciągłym funkcjonałem liniowym.
Dla każdej funkcji testowej $\varphi$ mamy
$\langle g(hT),\varphi\rangle=\langle hT,g\varphi\rangle
=\langle T,hg\varphi\rangle=\langle(gh)T,\varphi\rangle$
oraz $\langle1T,\varphi\rangle=\langle T,\varphi\rangle$.
Liniowość $M_g$ wynika z liniowości parowania względem $T$.

Jeśli $g$ zeruje się na otwartym $U\subseteq\Omega$, to
$g(e_U\varphi)=0$ dla $\varphi\in\mathcal D(U)$.
Jeśli natomiast $T|_U=0$, to
$g(e_U\varphi)=e_U((g|_U)\varphi)$, gdzie
$(g|_U)\varphi\in\mathcal D(U)$, zatem
$\langle gT,e_U\varphi\rangle=0$.
Stosując te obserwacje odpowiednio do
$U=\Omega\setminus\supp g$ i $U=\Omega\setminus\supp T$,
otrzymujemy $\supp(gT)\subseteq\supp g\cap\supp T$.
Ciągłość słaba-* wynika z tożsamości
$q_\varphi(gT)=q_{g\varphi}(T)$, a silna z
$q_B(gT)=q_{N_g(B)}(T)$, gdzie $B$ jest dowolnym zbiorem
ograniczonym w $\mathcal D(\Omega)$.
Zbiór $N_g(B)$ jest ograniczony, ponieważ $N_g$ jest
ciągły i liniowy (argument jak w Lemacie~\ref{lem:test-derivative20}).
Dla $j\in\{1,\ldots,d\}$ i $\alpha=e_j$ obliczamy
\[
 \begin{aligned}
 \langle\partial_j(gT),\varphi\rangle
 &=-\langle T,g\partial_j\varphi\rangle\\
 &=\langle\partial_jT,g\varphi\rangle
   +\langle T,(\partial_jg)\varphi\rangle,
 \end{aligned}
\]
czyli $\partial_j(gT)=g\partial_jT+(\partial_jg)T$.
Dla $\alpha=0$ wzór jest tożsamością $gT=gT$.
Załóżmy, że wzór zachodzi dla wszystkich wielowskaźników
rzędu $n$, i ustalmy $\gamma$ rzędu $n+1$.
Wybierzmy $j$ z $\gamma_j\ge1$ i połóżmy $\alpha=\gamma-e_j$.
Korzystając z przemienności pochodnych
(Twierdzenie~\ref{thm:dist-deriv-cont20}) oraz udowodnionej
reguły pierwszego rzędu dla każdego składnika, otrzymujemy
\[
\begin{aligned}
\partial^\gamma(gT)
={}&\sum_{\beta\le\alpha}\binom\alpha\beta
(\partial^{\beta+e_j}g)\partial^{\alpha-\beta}T\\
&+\sum_{\beta\le\alpha}\binom\alpha\beta
(\partial^\beta g)\partial^{\alpha-\beta+e_j}T.
\end{aligned}
\]
W pierwszej sumie podstawiamy $\nu=\beta+e_j$, a w drugiej
$\nu=\beta$. Przyjmując $\binom\alpha\mu=0$, gdy którakolwiek
współrzędna $\mu$ jest ujemna lub większa od odpowiedniej
współrzędnej $\alpha$, otrzymujemy
\[
\partial^\gamma(gT)=\sum_{\nu\le\gamma}
\left(\binom\alpha{\nu-e_j}+\binom\alpha\nu\right)
(\partial^\nu g)\partial^{\gamma-\nu}T.
\]
Ze zwykłej tożsamości Pascala dla $j$-tej współrzędnej,
po pomnożeniu przez $\prod_{k\ne j}\binom{\alpha_k}{\nu_k}$,
wynika $\binom\alpha{\nu-e_j}+\binom\alpha\nu=\binom\gamma\nu$.
Obejmuje to także skrajne wartości $\nu_j=0$ i $\nu_j=\gamma_j$.
Dowód indukcyjny jest zakończony.
\end{proof}

\begin{przyklad}[Mnożenie dystrybucji punktowych]
\label{ex:multiply-delta20}
Dla $a\in\R$ i $g\in C^\infty(\R;\K)$ mamy
\[
 g\delta_a=g(a)\delta_a,\qquad
 g\delta_a'=g(a)\delta_a'-g'(a)\delta_a.
\]
W szczególności $x\delta_0=0$ oraz $x\delta_0'=-\delta_0$.
Istotnie, dla $\varphi\in\mathcal D(\R)$ zachodzi
$\langle g\delta_a,\varphi\rangle=g(a)\varphi(a)$ oraz
\[
\begin{aligned}
\langle g\delta_a',\varphi\rangle
&=-(g\varphi)'(a)=-g'(a)\varphi(a)-g(a)\varphi'(a)\\
&=\langle g(a)\delta_a'-g'(a)\delta_a,\varphi\rangle.
\end{aligned}
\]
Podstawiając $g(x)=x$ i $a=0$, otrzymujemy obie
podane tożsamości szczególne.
\end{przyklad}

\begin{uwaga}[Dlaczego nie mnożymy dowolnych dystrybucji?]
\label{rem:no-product20}
Powyższa definicja wymaga \emph{gładkiego} mnożnika.
Nie ustanawia iloczynu dwóch dowolnych dystrybucji.
Przykładowo wyrażenie $H\delta_0$ nie ma kanonicznego
znaczenia na podstawie samej struktury $\mathcal D'(\R)$:
po zastąpieniu $\delta_0$ przez
$\rho_\varepsilon(x)=\varepsilon^{-1}\rho(x/\varepsilon)$,
gdzie $\rho\in C_c^\infty(\R;\R)$, $\rho\ge0$ oraz
$\int_\R\rho(x)\,\mathrm d\lambda_1(x)=1$,
otrzymujemy w sensie dystrybucji
\[
 H\rho_\varepsilon\longrightarrow c_\rho\delta_0,
 \qquad c_\rho=\int_0^\infty\rho(t)\,\mathrm d\lambda_1(t).
\]
Dla $\varepsilon>0$ iloczyn funkcji $H\rho_\varepsilon$
jest lokalnie całkowalny. Podstawienie $x=\varepsilon t$
daje dla każdej $\varphi\in\mathcal D(\R)$
\[
\langle T_{H\rho_\varepsilon},\varphi\rangle
=\int_0^\infty\rho(t)\varphi(\varepsilon t)\,\mathrm d\lambda_1(t).
\]
Całka dąży do $c_\rho\varphi(0)$ na mocy twierdzenia
o zbieżności zmajoryzowanej z majorantą
$\rho(t)p_0(\varphi)$.
Nawet silna zbieżność zachodzi: z
$|\varphi(\varepsilon t)-\varphi(0)|\le\varepsilon t p_1(\varphi)$
dla $t\ge0$ otrzymujemy
$q_B(T_{H\rho_\varepsilon}-c_\rho\delta_0)
\le\varepsilon\sup_{\varphi\in B}p_1(\varphi)
\int_0^\infty t\rho(t)\,\mathrm d\lambda_1(t)$
dla każdego niepustego ograniczonego $B\subseteq\mathcal D(\R)$.
Wszystkie takie $\rho_\varepsilon$ zbiegają do $\delta_0$
silnie, co wynika z analogicznego oszacowania całki po $\R$
z czynnikiem $|t|$. Wybierając znormalizowaną nieujemną
funkcję testową o nośniku w $(-1,0)$, otrzymujemy $c_\rho=0$,
a o nośniku w $(0,1)$ otrzymujemy $c_\rho=1$.
Istnienie obu funkcji wynika z Lematu~\ref{lem:cutoffs17}
i normalizacji przez dodatnią całkę.
Zatem takie regularizacje nie wyznaczają jednoznacznego
iloczynu $H\delta_0$; nie jest to zakaz konstruowania
iloczynów w dodatkowo określonych strukturach.
\end{uwaga}

\subsection{Operatory różniczkowe i ich transpozycja}

\begin{definicja}[Formalny operator transponowany]
\label{def:formal-transpose20}
Niech $m\in\N_0$ i $a_\alpha\in C^\infty(\Omega;\K)$
dla $|\alpha|\le m$. Dla funkcji gładkich i dystrybucji
rozważamy operator $P=\sum_{|\alpha|\le m}a_\alpha\partial^\alpha$.
Jego \emph{formalny operator transponowany} na $\mathcal D(\Omega)$
definiujemy wzorem
\[
 P^{\mathrm t}\varphi=
 \sum_{|\alpha|\le m}(-1)^{|\alpha|}
 \partial^\alpha(a_\alpha\varphi).
\]
Przy naszej konwencji parowania nie występuje sprzężenie
współczynników zespolonych. Nie należy zatem utożsamiać
formalnej transpozycji z formalnym sprzężeniem względem
zespolonego iloczynu skalarnego w $L^2$, ani z operatorem
sprzężonym określonym dla konkretnej dziedziny operatora.
W szczególnych przypadkach wzory mogą się pokrywać.
\end{definicja}

\begin{twierdzenie}[Transpozycja jako konstrukcja operatora na dystrybucjach]
\label{thm:transpose20}
Dla operatora $P$ z definicji~\ref{def:formal-transpose20}
oraz każdej dystrybucji $T\in\mathcal D'(\Omega;\K)$
zachodzą:
\begin{enumerate}
\item $P^{\mathrm t}:\mathcal D(\Omega)\to\mathcal D(\Omega)$
jest ciągły i liniowy;
\item wzór $\langle PT,\varphi\rangle=
\langle T,P^{\mathrm t}\varphi\rangle$ dla każdej
$\varphi\in\mathcal D(\Omega)$ wyznacza dystrybucję
zgodną z $\sum_{|\alpha|\le m}a_\alpha\partial^\alpha T$;
\item $P$ jest ciągły w topologiach słabej-* i silnej;
\item $\supp(PT)\subseteq\supp T$;
\item jeżeli $f\in C^m(\Omega)$, to $PT_f=T_{Pf}$.
\end{enumerate}
\end{twierdzenie}
\begin{proof}
Każdy składnik $P^{\mathrm t}$ jest równy
$(-1)^{|\alpha|}D_\alpha N_{a_\alpha}$, gdzie
$N_{a_\alpha}$ oznacza mnożenie funkcji testowych przez
$a_\alpha$. Obie operacje są ciągłe i liniowe na mocy
Lematu~\ref{lem:test-derivative20} i dowodu
Twierdzenia~\ref{thm:multiply20}. Suma jest skończona,
więc $P^{\mathrm t}$ jest ciągły i liniowy.
Ponieważ każda z operacji nie powiększa nośnika,
$\supp(P^{\mathrm t}\varphi)\subseteq\supp\varphi$.
Złożenie z $T$ jest ciągłym funkcjonałem liniowym.
Ponadto dla każdej funkcji testowej zachodzi
\[
\begin{aligned}
\left\langle\sum_{|\alpha|\le m}a_\alpha\partial^\alpha T,
\varphi\right\rangle
&=\sum_{|\alpha|\le m}\langle\partial^\alpha T,
a_\alpha\varphi\rangle\\
&=\sum_{|\alpha|\le m}(-1)^{|\alpha|}
\langle T,\partial^\alpha(a_\alpha\varphi)\rangle
=\langle T,P^{\mathrm t}\varphi\rangle.
\end{aligned}
\]
To dowodzi zgodności obu konstrukcji $PT$.
Zależności
$q_\varphi(PT)=q_{P^{\mathrm t}\varphi}(T)$ oraz
$q_B(PT)=q_{P^{\mathrm t}(B)}(T)$
dają obie ciągłości, ponieważ $P^{\mathrm t}(B)$ jest ograniczony.
Niech $U=\Omega\setminus\supp T$. Dla
$\varphi\in\mathcal D(U)$ funkcja $P^{\mathrm t}(e_U\varphi)$
ma nośnik zawarty w $\supp(e_U\varphi)\Subset U$,
jest więc przedłużeniem zerem funkcji testowej na $U$.
Wobec $T|_U=0$ mamy $\langle PT,e_U\varphi\rangle=0$,
a zatem $\supp(PT)\subseteq\supp T$.
Dla $f\in C^m(\Omega;\K)$, korzystając ze
Stwierdzenia~\ref{prop:classical-weak20} (dla $m=0$
pochodna zerowego rzędu jest identycznością), dostajemy
$\partial^\alpha T_f=T_{\partial^\alpha f}$.
Ponadto $a_\alpha T_{\partial^\alpha f}
=T_{a_\alpha\partial^\alpha f}$, gdyż obie strony mają
wartość $\int_\Omega a_\alpha(\partial^\alpha f)\varphi\,
\mathrm d\lambda_d$ na dowolnej funkcji testowej $\varphi$.
Sumowanie daje $PT_f=T_{Pf}$.
\end{proof}

\begin{przyklad}[Operator pierwszego rzędu]
\label{ex:first-order-P20}
W jednym wymiarze niech $a,b\in C^\infty(\R;\K)$
i $P=a(x)\frac{\mathrm d}{\mathrm dx}+b(x)$.
Z Definicji~\ref{def:formal-transpose20} otrzymujemy
$P^{\mathrm t}\varphi=-(a\varphi)'+b\varphi
=-a\varphi'+(b-a')\varphi$.
W $\R^d$ dla $b=(b_1,\ldots,b_d)\in C^\infty(\R^d;\K^d)$
i $c\in C^\infty(\R^d;\K)$ operator
$P=b\cdot\nabla+c=\sum_{j=1}^d b_j\partial_j+c$
ma formalną transpozycję
\[
P^{\mathrm t}\varphi=-\sum_{j=1}^d\partial_j(b_j\varphi)
+c\varphi=-b\cdot\nabla\varphi+(c-\operatorname{div}b)\varphi.
\]
Tutaj $\operatorname{div}b=\sum_{j=1}^d\partial_j b_j$,
a iloczyn $b\cdot\nabla$ nie zawiera sprzężenia zespolonego.
Druga równość wynika z klasycznej reguły Leibniza.
\end{przyklad}

\begin{wniosek}[Komutator pochodnej i mnożenia]
\label{cor:commutator20}
Dla $g\in C^\infty(\Omega)$ i $j\in\{1,\ldots,d\}$ mamy
$[\partial_j,M_g]=M_{\partial_jg}$ na $\mathcal D'(\Omega)$,
gdzie $[A,B]=AB-BA$ jest komutatorem operatorów.
\end{wniosek}
\begin{proof}
Jest to wzór Leibniza pierwszego rzędu:
$\partial_j(gT)-g\partial_jT=(\partial_jg)T$.
\end{proof}

\subsection{Równania różniczkowe w sensie dystrybucji}

\begin{twierdzenie}[Jądro pochodnej na przedziale]
\label{thm:zero-derivative20}
Niech $I\subseteq\R$ będzie niepustym przedziałem otwartym,
a $T\in\mathcal D'(I;\K)$. Wówczas $T'=0$
wtedy i tylko wtedy, gdy istnieje $c\in\K$ takie, że
$T=T_c$, czyli
$\langle T,\varphi\rangle=c\int_I\varphi(x)\,\mathrm d\lambda_1(x)$
dla każdej $\varphi\in\mathcal D(I)$.
\end{twierdzenie}
\begin{proof}
Wybierzmy niezerową nieujemną funkcję
$\eta_0\in C_c^\infty(I;\R)$ na mocy
Lematu~\ref{lem:cutoffs17}. Jej całka jest skończona
dzięki zwartemu nośnikowi i dodatnia, ponieważ $\eta_0$
jest dodatnia na pewnym niepustym otwartym podprzedziale.
Zatem
funkcja $\eta=\eta_0/\int_I\eta_0\,\mathrm d\lambda_1$
należy do $\mathcal D(I;\K)$ i ma całkę $1$.
Dla $\varphi\in\mathcal D(I)$ oznaczmy
$m_\varphi=\int_I\varphi(x)\,\mathrm d\lambda_1(x)$
i $h=\varphi-m_\varphi\eta$. Wtedy
$\int_I h\,\mathrm d\lambda_1=m_\varphi-m_\varphi=0$.
Przedłużamy $h$ zerem na $\R$ i określamy
$\psi(x)=\int_{-\infty}^xh(t)\,\mathrm d\lambda_1(t)$.
Przedłużenie jest gładkie, gdyż $\supp h\Subset I$.
Wybierzmy $u<v$ z $[u,v]\Subset I$ i
$\supp h\subseteq[u,v]$. Dla $x<u$ całka jest zerowa,
a dla $x>v$ jest równa całce po całej prostej, czyli również
zero. Z twierdzenia Newtona--Leibniza $\psi'=h$,
więc $\psi$ jest gładka i $\supp\psi\subseteq[u,v]$.
Jej ograniczenie do $I$ należy zatem do $\mathcal D(I)$. Jeśli $T'=0$, to
$0=\langle T',\psi\rangle=-\langle T,\psi'\rangle$,
więc $\langle T,\varphi\rangle=m_\varphi\langle T,\eta\rangle$.
Wybieramy $c=\langle T,\eta\rangle$; testowanie na $\eta$
pokazuje także jednoznaczność tej stałej.
Odwrotnie, dla $c\in\K$ i $\varphi\in\mathcal D(I)$
wybieramy przedział $[u,v]\Subset I$ zawierający
$\supp\varphi$ we wnętrzu. Wtedy
$\langle T_c',\varphi\rangle=-c\int_u^v\varphi'(x)\,
\mathrm d\lambda_1(x)=-c(\varphi(v)-\varphi(u))=0$.
\end{proof}

\begin{twierdzenie}[Istnienie pierwotnej dystrybucji]
\label{thm:primitive20}
Dla każdej dystrybucji $S\in\mathcal D'(\R;\K)$
istnieje $T\in\mathcal D'(\R;\K)$ takie, że $T'=S$.
Wszystkie takie pierwotne różnią się o stałą dystrybucję.
\end{twierdzenie}
\begin{proof}
Ustalmy $\eta\in\mathcal D(\R;\K)$ o całce $1$,
skonstruowaną jak w dowodzie
Twierdzenia~\ref{thm:zero-derivative20}, i
zdefiniujmy $m_\varphi=\int_\R\varphi(x)\,\mathrm d\lambda_1(x)$
oraz
\[
 R\varphi(x)=\int_{-\infty}^x
 \bigl(\varphi(t)-m_\varphi\eta(t)\bigr)\,\mathrm d\lambda_1(t).
\]
Całka funkcji $\varphi-m_\varphi\eta$ jest zerowa,
więc argument o nośniku z dowodu
Twierdzenia~\ref{thm:zero-derivative20} daje
$R\varphi\in\mathcal D(\R)$.
Liniowość $R$ wynika z liniowości całek.
Dla każdego zwartego $K\Subset\R$ wybieramy zwarty przedział
$L=[u,v]$ o dodatniej długości $\ell=v-u$, zawierający
$K\cup\supp\eta$. Wówczas $R(\mathcal D_K)\subseteq\mathcal D_L$.
Dla $\varphi\in\mathcal D_K$ mamy
$|m_\varphi|\le\lambda_1(K)p_0(\varphi)$.
Ponieważ $R\varphi$ znika poza $L$,
\[
p_0(R\varphi)\le\ell\bigl(p_0(\varphi)
+\lambda_1(K)p_0(\varphi)p_0(\eta)\bigr).
\]
Ustalmy $m\in\N_0$. Dla $1\le j\le m$ z twierdzenia Newtona--Leibniza
otrzymujemy $(R\varphi)^{(j)}=\varphi^{(j-1)}
-m_\varphi\eta^{(j-1)}$, zatem
$p_0((R\varphi)^{(j)})\le
(1+\lambda_1(K)p_m(\eta))p_m(\varphi)$.
Oba oszacowania dają dla każdego $m\in\N_0$
\[
p_m(R\varphi)\le
\max\{1,\ell\}(1+\lambda_1(K)p_m(\eta))p_m(\varphi).
\]
Ograniczenie $R|_{\mathcal D_K}:\mathcal D_K\to\mathcal D_L$
jest więc ciągłe, także gdy $m=0$ lub $K=\varnothing$.
Po złożeniu z ciągłym włożeniem $\mathcal D_L\to\mathcal D(\R)$
i zastosowaniu Twierdzenia~\ref{thm:LF-universal18}
otrzymujemy ciągłość $R:\mathcal D(\R)\to\mathcal D(\R)$.
Definiujemy
$\langle T,\varphi\rangle=-\langle S,R\varphi\rangle$.
Wtedy $T=-S\circ R$ jest ciągłym funkcjonałem liniowym,
czyli $T\in\mathcal D'(\R;\K)$.
Dla $\varphi\in\mathcal D(\R)$ zachodzi
$m_{\varphi'}=\int_\R\varphi'\,\mathrm d\lambda_1=0$,
a z twierdzenia Newtona--Leibniza
$R(\varphi')(x)=\int_{-\infty}^x\varphi'(t)\,
\mathrm d\lambda_1(t)=\varphi(x)$, gdyż $\varphi$ znika
dla dostatecznie małych argumentów.
Stąd
$\langle T',\varphi\rangle=-\langle T,\varphi'\rangle
=\langle S,R\varphi'\rangle=\langle S,\varphi\rangle$.
Jeżeli $T_1'=T_2'=S$, to $(T_1-T_2)'=0$, a z
Twierdzenia~\ref{thm:zero-derivative20} wynika $T_1-T_2=T_c$
dla pewnego $c\in\K$. Odwrotnie, dodanie $T_c$ do
dowolnej pierwotnej nie zmienia jej pochodnej.
\end{proof}

\begin{wniosek}[Równanie pierwszego rzędu]
\label{cor:ode20}
Niech $a\in C^\infty(\R;\K)$ i
$S\in\mathcal D'(\R;\K)$. Wszystkie dystrybucyjne
rozwiązania równania $T'+aT=S$ mają postać
$T=e^{-A}V$, gdzie $A(x)=\int_0^xa(t)\,\mathrm d\lambda_1(t)$,
a $V$ jest dowolną pierwotną dystrybucją $e^AS$.
Dwa rozwiązania różnią się o dystrybucję regularną
$T_{c e^{-A}}$ dla pewnego $c\in\K$.
\end{wniosek}
\begin{proof}
Z twierdzenia Newtona--Leibniza $A'=a$, a funkcje $e^A$
i $e^{-A}$ są gładkie, także dla zespolonego $a$.
Z Twierdzenia~\ref{thm:multiply20} operatory mnożenia
przez te funkcje są wzajemnie odwrotne.
Reguła Leibniza daje dla każdej dystrybucji $T$
$(e^AT)'=e^AT'+(e^A)'T=e^A(T'+aT)$.
Zatem $T'+aT=S$ wtedy i tylko wtedy, gdy
$(e^AT)'=e^AS$. Z Twierdzenia~\ref{thm:primitive20}
istnieje pierwotna $V$ dystrybucji $e^AS$ i wszystkie
pozostałe mają postać $V+T_c$, $c\in\K$.
Mnożąc przez $e^{-A}$, otrzymujemy wszystkie rozwiązania
$e^{-A}V+T_{ce^{-A}}$.
\end{proof}

\begin{przyklad}[Rozwiązania dystrybucyjne i skoki]
\label{ex:forced-ode20}
Dla $\lambda\in\K$ rozwiązujemy równanie
$T'+\lambda T=\delta_0$ na $\R$.
Funkcja $f(x)=e^{-\lambda x}H(x)$ jest lokalnie całkowalna,
jest klasy $C^1$ poza zerem, ma granice $f(0-)=0$ i
$f(0+)=1$, a jej klasyczna pochodna poza zerem jest równa
$-\lambda f$ i jest lokalnie całkowalna.
Z Twierdzenia~\ref{thm:jump20} dla $E_\lambda=T_f$ mamy
$E_\lambda'=-\lambda E_\lambda+\delta_0$, czyli
$E_\lambda'+\lambda E_\lambda=\delta_0$.
Z Wniosku~\ref{cor:ode20}, przy $A(x)=\lambda x$,
wszystkie rozwiązania mają więc postać
$T=E_\lambda+cT_{e^{-\lambda x}}$, $c\in\K$.
Dla $\lambda=0$ otrzymujemy równanie $T'=\delta_0$
i pierwotną $T_H$.
Ponadto dla $x_+(x)=xH(x)$ z reguły Leibniza mamy
$(T_{x_+})'=T_H+x\delta_0=T_H$, ponieważ
$x\delta_0=0$ (Przykład~\ref{ex:multiply-delta20}).
Różniczkując ponownie, dostajemy
$\frac{\mathrm d^2}{\mathrm dx^2}T_{x_+}=\delta_0$.
\end{przyklad}

\begin{twierdzenie}[Znikanie gradientu na zbiorze spójnym]
\label{thm:zero-gradient20}
Niech $\Omega\subseteq\R^d$ będzie niepustym spójnym
zbiorem otwartym i $T\in\mathcal D'(\Omega;\K)$.
Jeżeli $\partial_jT=0$ dla każdego $j=1,\ldots,d$,
to istnieje $c\in\K$ takie, że
$\langle T,\varphi\rangle=c\int_\Omega\varphi(x)\,
\mathrm d\lambda_d(x)$ dla każdej $\varphi\in\mathcal D(\Omega)$.
\end{twierdzenie}
\begin{proof}
Najpierw rozważmy otwarty prostokąt
$Q=I_1\times\cdots\times I_d$ o domknięciu zwartym w $\Omega$.
Przez $T_Q=T|_Q$ oznaczmy ograniczenie dystrybucji do $Q$.
Zgodność różniczkowania z przedłużaniem zerem, wykazana
w dowodzie Twierdzenia~\ref{thm:dist-deriv-cont20},
daje $\partial_jT_Q=0$ dla każdego $j$.
Dla każdego $j$ wybieramy $\eta_j\in\mathcal D(I_j)$
o całce $1$ jak w dowodzie Twierdzenia~\ref{thm:zero-derivative20}
i dla $\varphi\in\mathcal D(Q)$ określamy operator
\[
 (P_j\varphi)(x)=\eta_j(x_j)
 \int_{I_j}\varphi(x_1,\ldots,x_{j-1},t,x_{j+1},\ldots,x_d)
 \mathrm d\lambda_1(t).
\]
Dla $\varphi\in\mathcal D(Q)$ wybieramy zwarte przedziały
$L_k\Subset I_k$ zawierające rzuty $\supp\varphi$
oraz $\supp\eta_k$ dla $k=1,\ldots,d$.
Całkowanie względem $t$ w definicji $P_j$ można ograniczyć
do $L_j$. Różniczkowanie pod całką względem pozostałych
zmiennych jest dozwolone, gdyż całkujemy funkcję gładką
po ustalonym zwartym przedziale; lokalnie wszystkie
pochodne są ograniczone na odpowiednich iloczynach zwartych.
Zatem $P_j\varphi$ jest gładka.
Jeżeli $x_j\notin L_j$, znika czynnik $\eta_j(x_j)$,
a jeżeli $x_k\notin L_k$ dla pewnego $k\ne j$, znika cała
funkcja podcałkowa. Stąd $\supp(P_j\varphi)\subseteq
L_1\times\cdots\times L_d$, czyli $P_j\varphi\in\mathcal D(Q)$.
Liniowość $P_j$ wynika z liniowości całki.
Dla $j\ne k$ oba złożenia $P_jP_k$ i $P_kP_j$ dają
czynnik $\eta_j(x_j)\eta_k(x_k)$ pomnożony przez całkę
podwójną funkcji $\varphi$ po $I_j\times I_k$ przy pozostałych
zmiennych ustalonych. Twierdzenie Fubiniego (całka jest
bezwzględnie zbieżna) daje ich równość.

Dla $\psi\in\mathcal D(Q)$ oznaczmy
$h_j=(\operatorname{Id}-P_j)\psi$.
Z definicji $P_j$ i $\int_{I_j}\eta_j\,\mathrm d\lambda_1=1$
wynika, że całka $h_j$ po $x_j$ jest zerowa dla każdego
ustalenia pozostałych zmiennych.
Przedłużamy $h_j$ zerem na $\R^d$ i określamy
\[
\theta_j(x)=\int_{-\infty}^{x_j}
h_j(x_1,\ldots,x_{j-1},t,x_{j+1},\ldots,x_d)\,
\mathrm d\lambda_1(t).
\]
Wybieramy zwarte przedziały $L_k=[u_k,v_k]\Subset I_k$
tak, aby ich iloczyn zawierał $\supp h_j$.
Całka w definicji $\theta_j$ jest zerowa, gdy
$x_j<u_j$ albo $x_j>v_j$ (w drugim przypadku dzięki
zerowej całce po całej zmiennej $x_j$), a także gdy
$x_k\notin L_k$ dla pewnego $k\ne j$.
Zatem $\supp\theta_j\subseteq L_1\times\cdots\times L_d$.
Z twierdzenia Newtona--Leibniza i różniczkowania pod całką
względem pozostałych zmiennych $\theta_j$ jest gładka i
$\partial_j\theta_j=h_j$. Ograniczenie $\theta_j$ do $Q$
jest funkcją testową. Ponieważ $\partial_jT_Q=0$,
\[
\langle T_Q,(\operatorname{Id}-P_j)\psi\rangle
=\langle T_Q,\partial_j\theta_j\rangle
=-\langle\partial_jT_Q,\theta_j\rangle=0.
\]
Ostatnia równość mówi, że
$\langle T_Q,\psi\rangle=\langle T_Q,P_j\psi\rangle$
dla każdej $\psi\in\mathcal D(Q)$ i każdego $j$.
Stosując ją kolejno do $\varphi,P_1\varphi,\ldots,
P_{d-1}\cdots P_1\varphi$ i korzystając z przemienności
operatorów, otrzymujemy
$\langle T_Q,\varphi\rangle=\langle T_Q,P_1\cdots P_d\varphi\rangle$.
Wielokrotne zastosowanie definicji $P_j$ i twierdzenia
Fubiniego daje
$P_1\cdots P_d\varphi=(\int_Q\varphi\,\mathrm d\lambda_d)
\eta_1\cdots\eta_d$, gdzie iloczyn $\eta_1\cdots\eta_d$
oznacza funkcję $x\mapsto\prod_{j=1}^d\eta_j(x_j)$.
Ma ona zwarty nośnik w $Q$ i całkę $1$.
Zatem $T_Q=T_{c_Q}$, gdzie
$c_Q=\langle T_Q,\eta_1\cdots\eta_d\rangle$.
Jeśli $Q\cap Q'\ne\varnothing$, wybieramy funkcję testową
$\zeta$ o całce $1$ i nośniku w $Q\cap Q'$ (lemat o
funkcjach odcięcia i normalizacja). Testując oba ograniczenia
na przedłużeniach $\zeta$ zerem, otrzymujemy
$c_Q=\langle T,\zeta\rangle=c_{Q'}$, gdzie w środkowym
wyrażeniu $\zeta$ jest przedłużona zerem na $\Omega$.

Każdy punkt $x\in\Omega$ należy do otwartego prostokąta
$Q$ o domknięciu zwartym w $\Omega$. Przypisujemy mu
$c(x)=c_Q$; poprzedni argument pokazuje, że wartość nie
zależy od wyboru prostokąta. Funkcja $c(x)$ jest lokalnie
stała, gdyż na każdym takim $Q$ jest równa $c_Q$.
Ustalmy $x_0\in\Omega$. Zbiór
$\{x\in\Omega:c(x)=c(x_0)\}$ i jego dopełnienie są otwarte.
Pierwszy zbiór jest niepusty, więc ze spójności $\Omega$
jest równy $\Omega$. Oznaczmy wspólną stałą przez $c$.
Dystrybucja $T-T_c$ zeruje się na każdym takim prostokącie,
a prostokąty pokrywają $\Omega$. Z definicji nośnika
$\supp(T-T_c)=\varnothing$, a z Lematu~\ref{lem:distribution-local19}
wynika $T-T_c=0$. To daje żądaną równość na całym $\Omega$.
Stała $c$ jest jednoznaczna, gdyż można wybrać funkcję
testową o całce $1$ w dowolnym prostokącie $Q\Subset\Omega$
i obliczyć na niej wartość $T$.
\end{proof}

\begin{uwaga}[Granice twierdzeń]
\label{rem:limits20}
Niech $d\ge2$, $I\subseteq\R$ będzie niepustym przedziałem
otwartym, a $V\subseteq\R^{d-1}$ niepustym zbiorem otwartym.
Niech $T\in\mathcal D'(I\times V;\K)$.
Warunek $\partial_1T=0$ oznacza, że $T$ zależy od funkcji
testowej wyłącznie przez jej całkę po pierwszej współrzędnej.
Istotnie, wybierzmy $\eta\in\mathcal D(I)$ o całce $1$.
Dla $\varphi\in\mathcal D(I\times V)$ połóżmy
$A\varphi(y)=\int_I\varphi(t,y)\,\mathrm d\lambda_1(t)$.
Nośnik $A\varphi$ mieści się w zwartym rzucie
$\supp\varphi$ na $V$, a gładkość wynika z różniczkowania
pod całką po ustalonym zwartym podprzedziale $I$,
zawierającym rzut nośnika na $I$.
Funkcja $h(t,y)=\varphi(t,y)-\eta(t)A\varphi(y)$ ma zerową
całkę po $I$ dla każdego $y$.
Po przedłużeniu zerem jej pierwotna względem $t$,
$\theta(t,y)=\int_{-\infty}^t h(s,y)\,\mathrm d\lambda_1(s)$,
jest gładka. Jej nośnik jest zawarty w $[u,v]\times K$,
gdzie $[u,v]\Subset I$ zawiera rzut $\supp h$ na $I$,
a $K\Subset V$ jest rzutem $\supp h$ na $V$.
Zatem $\theta|_{I\times V}$ jest funkcją testową i
$\langle T,h\rangle=-\langle\partial_1T,\theta\rangle=0$.
Otrzymujemy
\[
\langle T,\varphi\rangle
=\left\langle S,\int_I\varphi(t,\cdot)\,\mathrm d\lambda_1(t)\right\rangle,
\qquad
\langle S,\psi\rangle=\langle T,\eta\otimes\psi\rangle,
\]
gdzie $(\eta\otimes\psi)(t,y)=\eta(t)\psi(y)$.
Odwzorowanie $\psi\mapsto\eta\otimes\psi$ jest ciągłe
na każdym etapie $\mathcal D_K(V)$: nośnik obrazu mieści
się w $\supp\eta\times K$, a jego seminorma $p_m$ nie
przekracza $p_m(\eta)p_m(\psi)$. Z własności uniwersalnej
granicy induktywnej jest ciągłe na $\mathcal D(V)$.
Stąd $S\in\mathcal D'(V;\K)$. Dystrybucja $S$ jest jednoznacznie wyznaczona,
ponieważ $A(\eta\otimes\psi)=\psi$ dla każdej $\psi\in\mathcal D(V)$.
Odwrotnie, gdy $T$ ma podaną reprezentację, to
$A(\partial_1\varphi)=0$ z twierdzenia Newtona--Leibniza
i zwartego nośnika $\varphi$, więc $\partial_1T=0$.
Dla przykładu regularna dystrybucja odpowiadająca funkcji
$(t,y)\mapsto y_1$ ma zerową pierwszą pochodną, ale
nie jest stała na niepustym zbiorze otwartym $I\times V$.
Dla układu $\partial_jT=S_j$, gdzie
$T,S_1,\ldots,S_d\in\mathcal D'(\Omega;\K)$, konieczne są warunki zgodności
$\partial_kS_j=\partial_jS_k$, wynikające z przemienności
pochodnych dystrybucyjnych. Ich wystarczalność globalna
może zależeć od topologii zbioru otwartego; nie wynika z tego
rozdziału.
\end{uwaga}

\begin{uwaga}[Związek z operatorami na TVS]
\label{rem:tvs-derivative20}
Konstrukcja $T\mapsto PT$ nie wymaga normy na
$\mathcal D(\Omega)$ ani na $\mathcal D'(\Omega)$.
Opiera się na ciągłości operatora $P^{\mathrm t}$ na
lokalnie wypukłej przestrzeni funkcji testowych.
W ten sposób pochodna dystrybucyjna jest szczególnym
przypadkiem transpozycji ciągłego operatora liniowego
między przestrzeniami liniowo-topologicznymi.
\end{uwaga}
